(1) 真空中一均匀线性各向同性材料椭球的三个半轴为 \(a,b,c>0\)，相对介电常数 \(\varepsilon_r>1\)，置于均匀静外电场 \(\boldsymbol E_0\)。已知椭球均匀极化产生内部退极化场 \(\boldsymbol E_{\rm dep}=-\boldsymbol L\boldsymbol P/\varepsilon_0\)，在主轴基下 \(\boldsymbol L=\mathrm{diag}(L_x,L_y,L_z)\)，\(L_i>0,\sum L_i=1\)。求内部电场、极化与极化率张量。
(2) 由 \(L_i=(abc/2)\int_0^\infty ds/[(s+a_i^2)\sqrt{(s+a^2)(s+b^2)(s+c^2)}]\) 验证球形 \(a=b=c=R\) 时 \(L_i=1/3\)，得到经典介电球极化率，并求 \(\varepsilon_r\to\infty\) 导体极限。
(3) 在外场 \(\boldsymbol E_0\) 与两主轴构成的平面中，设极化率分别为 \(\alpha_x,\alpha_y\)，外场与第一主轴夹角 \(\theta\)。求固定外场时的可逆诱导偶极能量及绕垂直轴的力矩，判断 \(\alpha_x>\alpha_y\) 时稳定取向。
(4) 考虑长旋转椭球 \(a>b=c\)，离心率 \(e=\sqrt{1-b^2/a^2}\)，已知 \(L_x=(1-e^2)[\ln((1+e)/(1-e))-2e]/(2e^3)\)。求 \(e\to0\) 与 \(e\to1^-\) 的极限，解释长轴增强极化率和线性介质模型失效条件。
